% SHARD-ID: AQM-87-QSA-OPEN-PROBLEMS-NEXT-TARGETS
% SHARD-TITLE: Open Problems and Next Catalogue Targets
% SHARD-SUMMARY: Closes the QSA meta-plan with a prioritized gap ledger rather than new derivations.
% SHARD-SUMMARY: Turns the preserved Step 21--22 blockers into explicit source, convention, computation, and run targets.
% SHARD-SUMMARY: Records parafermion, fusion-category, finite-ring, completion, dynamics, and noncommutative/nonunital boundaries as questions or proposals.
% SHARD-KEYWORDS: quantum-system association, open problems, next targets, blockers, finite rings, fusion category, completions, dynamics

\section{Open Problems and Next Catalogue Targets}
\label{sec:qsa-open-problems-next-targets}

\subsection{Status and Local Anchors}

\textbf{Status: New gap/synthesis shard.}  This is QSA Step 23 of the plan
in \path{docs/quantum_system_associations/}.  It adds no source, convention
entry, script, run bundle, generated data, Beads update, database row,
representation classification, dynamics principle, or theorem.  It closes
the first QSA meta-plan by recording the next work that would be needed
before further catalogue rows become evidence.

The governing vocabulary is AQM-65 and \texttt{CONVENTIONS.md} convention
\texttt{(ap)}.  The table tokens are those of convention \texttt{(au)}.  The
local QSA anchors are AQM-72 and AQM-73 for parafermion and fusion-category
gaps, AQM-78 for finite Weyl-Heisenberg representation boundaries, AQM-80
for finite-ring database status, AQM-84 for infinite-ring and completion
boundaries, AQM-85 for the agreement/inequivalence ledger, and AQM-86 for
degenerate and over-restrictive cases.

The index-level anchors are the source-topic statuses in \texttt{INDEX.md}
and the finite-ring database run from 2026-05-26.  Those anchors say what
exists locally; they are not imported as new mathematical sources in this
shard.

\subsection{Reading Rule}

Every item below is an open question or a proposal.  A proposed target is
not a claim that the target is correct, natural, or unique.  It is a record
of missing evidence needed before a later shard may add a result row.

The word \texttt{blocked} keeps the AQM-85--AQM-86 meaning: a source,
convention, computation, helper, audit, or run is missing.  It is not a
mathematical obstruction.  Conversely, an \texttt{available} local anchor is
not a final equivalence statement unless AQM-85's evidence rule supplies a
specific object, workflows, choices, and a local map, invariant, or
obstruction.

\subsection{Prioritized Dependency Ledger}

\begin{center}
\scriptsize\setlength{\tabcolsep}{3pt}
\begin{tabular}{p{0.07\linewidth}p{0.16\linewidth}p{0.28\linewidth}p{0.31\linewidth}p{0.10\linewidth}}
\toprule
Pri. & Bucket & Open question or proposal & Next target must supply & Anchors \\
\midrule
1 &
Source and convention acquisition &
Question: which parafermion or fusion-category association should be the
first technical extension beyond CCR, CAR, Grassmann-Weyl, and finite
AND/OR grammar? &
Acquire and register the missing parafermion or fusion-category source;
record conventions for mode algebra or categorical data, including hom
spaces, simple labels, tensor rules, associator data, pivotal/spherical data,
and any lattice input before using them. &
AQM-72, AQM-73 \\
\addlinespace
2 &
Finite-ring database, residue, and cotangent audits &
Question: which standard finite-ring rows can become database-backed rather
than local exceptions or review-export notes? &
Produce a populated and audited finite-ring run; record deterministic
residue-site ordering and database-backed intrinsic/relative cotangent
computations before extending the AQM-80 matrix. &
AQM-80, AQM-81 \\
\addlinespace
3 &
Quotient collapse and two-direction Artin rows &
Question: which quotient presentations collapse to earlier examples, and
which two-direction Artin rings support a checked thickened Weyl layer? &
For quotient collapse, add a local derivation or certificate beyond the
\(\mathbb Z/6\mathbb Z\) anchor.  For two-direction Artin examples, add the
missing source or exact derivation, generating-character convention, and
checked computation. &
AQM-80, AQM-82, AQM-86 \\
\addlinespace
4 &
Finite Weyl and projective representation boundaries &
Question: for which finite rings or finite abelian-group pairings can the
Weyl-Heisenberg/projective-representation step be classified beyond the
locally sourced finite-field, finite Frobenius-ring Pauli, and Solomon-style
generalized-Heisenberg slots? &
Name the finite label data, central character, commutator pairing, and
source or derivation.  If the claim is computational, add a run checking the
projective product law, commutator law, irreducibility or decomposition
invariant, and operator-basis count. &
AQM-31, AQM-78 \\
\addlinespace
5 &
Single-particle, Fock, and statistics comparisons &
Question: when do the AQM-74 one-particle maps extend to symmetric,
antisymmetric, para, or full Fock sectors, and when do they fail? &
Record the statistics convention, particle-number grading, completion or
finite-sector cutoff, and an explicit comparison map or obstruction.  For
para-statistics, first acquire a source and mode-algebra convention. &
AQM-70, AQM-72, AQM-74 \\
\addlinespace
6 &
Infinite, generic, analytic, and completion boundaries &
Question: which topology or completion should turn finite-support closed
point nets, generic sectors, or formal infinite-volume profiles into
analytic quantum systems? &
Choose and source the topology, measure or character, Pontryagin-dual data,
norm or \(C^*\)-completion, adelic or profinite restricted product, and
representation theorem before claiming a Hilbert-space model. &
AQM-84 \\
\addlinespace
7 &
Dynamics selection &
Question: what principle, if any, should select Hamiltonians, time
evolutions, actions, or Witten-index-style data after a kinematical
association has been fixed? &
Add a dynamics convention naming the sign, domain, adjoint, grading,
boundary/support condition, and source or checked invariant.  Until then,
Hamiltonians and time evolution remain outside QSA kinematics. &
AQM-65, AQM-78 \\
\addlinespace
8 &
Noncommutative or nonunital extension boundary &
Question: what should replace \(\Spec R\), residue fields, cotangent data,
and tensor-site assembly if the input ring is noncommutative or lacks a
unit? &
Acquire sources and record conventions for the chosen replacement of points
and local data before adding examples.  Possible prompts such as
unitization, idempotent/corner data, or module-category input remain only
prompts until sourced. &
AQM-75--AQM-83 \\
\bottomrule
\end{tabular}
\end{center}

\subsection{Blockers Carried Forward}

The blockers below are carried forward without reinterpretation.
\[
  \begin{gathered}
  \texttt{aqm-qsa-src02},\quad
  \texttt{aqm-qsa-frres01},\quad
  \texttt{aqm-qsa-frcot01},\\
  \texttt{aqm-qsa-collapse01},\quad
  \texttt{aqm-qsa-artin01}.
  \end{gathered}
\]

\begin{description}
\item[\texttt{aqm-qsa-src02}.]
Question: which missing sources must be acquired before parafermions,
fusion categories, Levin-Wen/string-net data, or related categorical
association rows become technical report content?  The next target must name
the source topic, local manifest path, retrieval route, source locators, and
the conventions that the source fixes.
\item[\texttt{aqm-qsa-frres01}.]
Question: can the finite-ring database supply certified maximal/residue
decompositions and deterministic site orderings for the standard finite-ring
rows?  The next target must supply a populated run, audit result, and report
or schema hooks before using those rows as database-backed evidence.
\item[\texttt{aqm-qsa-frcot01}.]
Question: can intrinsic and relative cotangent rows be computed and audited
from the finite-ring database rather than by local hand derivations?  The
next target must name the base ring convention, cotangent algorithm,
verification invariant, and run artifact.
\item[\texttt{aqm-qsa-collapse01}.]
Question: which quotient presentations are checked collapses rather than
new examples?  The next target must provide a local derivation or certificate
for each claimed collapse beyond the \(\mathbb Z/6\mathbb Z\) anchor.
\item[\texttt{aqm-qsa-artin01}.]
Question: which two-direction or arbitrary Artin/Frobenius examples have the
generating-character and thickened Weyl data needed for a nilpotent-sensitive
row?  The next target must supply the source or exact derivation, the
character convention, and a checked example.
\end{description}

\subsection{Catalogue Hygiene Rules for the Next Round}

\begin{itemize}
\item Proposal: every new row should begin by naming the object, the selected
  workflow, the extra choices, and the intended output before any comparison
  word is used.
\item Proposal: every comparison target should declare whether it needs a
  missing source, convention, computation, or run; otherwise its status should
  remain \texttt{blocked} or \texttt{unknown}.
\item Proposal: every finite-ring extension should state whether it is a
  residue-site row, cotangent-first row, reduced field-label row,
  nilpotent-sensitive row, quotient-collapse row, or database-audit row.
\item Proposal: every infinite or analytic extension should state whether it
  is finite-support, algebraic quasi-local, incomplete tensor, generic-sector,
  formal infinite-volume, or completion status in the AQM-84 vocabulary.
\item Proposal: dynamics should stay out of the QSA catalogue unless a later
  shard explicitly chooses and sources a dynamics-selection principle.
\end{itemize}

\subsection{Close of the First QSA Pass}

The first QSA pass has produced a controlled vocabulary, a standard test
object list, finite-set workflows, finite symplectic field-label workflows,
boson/fermion review rows, geometric point and cotangent workflows,
finite-ring residue/nilpotent/product rows, infinite boundary statuses,
an agreement/inequivalence ledger, and a failure-mode catalogue.  This shard
does not promote that material to a complete classification.  It records the
evidence still missing before the next catalogue pass can responsibly expand
the association zoo.
